% 赛后期刊论文骨架（LaTeX）。投稿前按目标期刊模板替换 documentclass。
\documentclass[review]{elsarticle} % 以 Elsevier 为例；MDPI 用 mdpi 模板
\usepackage{graphicx,amsmath,amssymb,booktabs,algorithm,algpseudocode}

\title{Multi-Agent Simulation and <Learning/Optimization> for <Scenario>: A <Contribution> Approach}

\author[1]{First Author\corref{cor1}}
\cortext[cor1]{Corresponding author}
\address[1]{School of ..., University ..., City, China}

\begin{abstract}
    <Background>... <Method>... <Key results with numbers>... <Significance>...
\end{abstract}

\begin{keyword}
multi-agent simulation \sep reinforcement learning \sep scheduling \sep calibration
\end{keyword}

\begin{document}
\maketitle

\section{Introduction}\label{sec:intro}
    Motivation, gap, contributions (3--4 bullets), paper organization.

\section{Related Work}\label{sec:related}
    \subsection{Agent-based / discrete-event simulation}
    \subsection{Multi-agent reinforcement learning / optimization}
    \subsection{Simulation calibration}

\section{Problem Formulation}\label{sec:problem}
    Notation, objective, constraints, KPI definitions.

\section{Method}\label{sec:method}
    \subsection{Simulation model} (ODD-style description)
    \subsection{Learning / optimization algorithm}
    \subsection{Calibration and training details}

\section{Experiments}\label{sec:exp}
    \subsection{Setup} (data, metrics, baselines, hardware, seeds)
    \subsection{Main results} (tables, figures)
    \subsection{Ablation study}
    \subsection{Sensitivity and statistical significance}

\section{Conclusion}\label{sec:conclusion}
    Contributions, limitations, future work.

\bibliographystyle{elsarticle-num}
\bibliography{references}
\end{document}
